\documentclass[11pt]{article}
\usepackage[margin=0.9in]{geometry}
\usepackage{xcolor}
\usepackage{booktabs}
\usepackage{tabularx}
\usepackage{enumitem}
\usepackage{titlesec}
\usepackage[hidelinks]{hyperref}

\definecolor{accent}{HTML}{315B7D}
\titleformat{\section}{\large\bfseries\color{accent}}{\thesection.}{0.5em}{}
\titlespacing*{\section}{0pt}{1.0em}{0.35em}
\setlength{\parindent}{0pt}
\setlength{\parskip}{0.45em}

\begin{document}
\begin{center}
{\LARGE\bfseries Stage 1 Registered Report}\\[0.25em]
{\large Sleep restriction and confidence calibration in medical decisions}\\[0.55em]
{\small Priya Nadeau, Samuel Okoro, and Lena Fischer}\\
{\small Protocol version 1.2 \quad Registered 24 July 2026}
\end{center}
\color{accent}\rule{\textwidth}{1pt}\color{black}

\section{Research question and hypotheses}
Does one night of moderate sleep restriction impair clinicians' ability to
calibrate diagnostic confidence? The primary hypothesis is that restriction
increases Brier score relative to a rested control condition. Accuracy,
response time, and metacognitive efficiency are secondary outcomes.

\section{Design}
The study uses a randomized, counterbalanced crossover design with a seven-day
washout. The confirmatory sample is 96 licensed clinicians recruited before
unblinding. The protocol, materials, analysis code, and simulated-data checks
are archived with immutable version identifiers.

\begin{tabularx}{\textwidth}{@{}lX@{}}
\toprule
\textbf{Decision} & \textbf{Pre-specified rule} \\
\midrule
Primary outcome & Mean Brier score across 40 diagnostic vignettes \\
Exclusions & Failed attention check or protocol deviation recorded before scoring \\
Missing data & Hierarchical model using all observed trials, with sensitivity analysis \\
Stopping rule & Recruitment ends at 96 complete crossover participants \\
Alpha and tails & Two-sided $\alpha=0.05$ for the single primary contrast \\
\bottomrule
\end{tabularx}

\section{Sampling and power}
Simulation under the planned hierarchical model gives 90\% power for a
standardized within-participant effect of 0.32. The simulation script varies
trial correlation, attrition, and heteroscedasticity. No interim outcome test
will be performed.

\section{Analysis plan}
The primary model includes condition as a fixed effect and participant and
vignette as varying intercepts. We will report the estimated contrast, 95\%
interval, standardized effect, model diagnostics, and the full multiverse
defined below.

\begin{enumerate}[leftmargin=1.5em,itemsep=0.15em]
\item Fit the confirmatory model to the locked analysis dataset.
\item Run prior predictive and posterior predictive checks.
\item Repeat with the pre-specified complete-case and robust-likelihood models.
\item Label every non-registered analysis as exploratory.
\end{enumerate}

\section{Interpretation table}
\begin{tabularx}{\textwidth}{@{}lX@{}}
\toprule
\textbf{Outcome} & \textbf{Planned interpretation} \\
\midrule
Positive primary effect & Evidence that restriction worsens calibration \\
Interval excludes target effect & Evidence against a practically meaningful effect \\
Inconclusive interval & Insufficient precision, without a directional claim \\
\bottomrule
\end{tabularx}

\section{Ethics, data, and deviations}
Insert the ethics approval or exemption, consent process, compensation, privacy
controls, and retention period. Stage 2 will include a dated deviation log,
including deviations with no effect on the conclusions. De-identified data and
reproduction code will be released under the licenses stated in the repository.

\end{document}
